\section{Introduction}
\label{sec:intro}

Autonomous agents built on \LLMs now select and invoke tools---calculators, search
engines, code interpreters, translation services---to answer questions that exceed their
parametric knowledge.  A single wrong tool call cascades: it produces inputs unsuited to
the downstream task, forcing re-planning or causing outright failure.  As agent
deployments reach high-stakes settings in finance, healthcare, and software engineering,
understanding \emph{why} a model picks one tool over another is no longer optional.

\para{What do we currently know?}  Behavioral benchmarks such as \toolbench
\citep{qin2023toolllm} and Berkeley Function Calling Leaderboard \citep{berkeleyfcl2023}
measure selection accuracy, but treat the model as a black box.  \citet{blankenstein2025biasbust}
show that semantic alignment between a query and a tool's description is the strongest
behavioral predictor of selection, and that positional biases also operate---yet neither
work explains the internal computation that produces these patterns.  Closest to our
goals, \citet{healy2026internal} demonstrate that the \emph{final-layer} hidden state of
a model encodes information sufficient to detect hallucinated tool calls, confirming that
internal representations carry tool-selection-relevant signal.  However, no prior work
identifies \emph{which layers} perform the computation, \emph{whether} it differs from
coarse category routing, or \emph{how much} semantic content versus positional order
drives the decision.

\para{Our approach.}  We address these gaps with four empirical experiments designed to
characterize tool selection from behavioral and mechanistic angles
(see \figref{fig:overview}).  \Figref{fig:overview} provides an overview of all key
metrics.

\begin{figure}[t]
    \centering
    \includegraphics[width=0.95\linewidth]{fig7_summary.png}
    \caption{Summary of all experimental results.  Semantic similarity achieves 40\%
             top-1 accuracy (vs.\ 25\% chance), selection is stable across 90\% of
             orderings, and the correct tool position is decodable from layer 9 with
             68\% probe accuracy---all pointing to a semantic-content-driven but
             multi-layer computation.}
    \label{fig:overview}
\end{figure}

\begin{enumerate}[leftmargin=*,itemsep=0pt,topsep=2pt]
    \item \textbf{Semantic similarity analysis}: We use \minilm sentence embeddings to
          measure how well cosine similarity between a query and each tool's description
          predicts actual \gptmini selections.
    \item \textbf{Positional bias testing}: We rotate the order of tool candidates across
          four cyclic shifts and measure whether the model's selection follows semantic
          content or list position.
    \item \textbf{Causal description perturbation}: We swap tool descriptions
          adversarially and test whether the model's selection tracks the description
          (semantics) or the tool name (identity).
    \item \textbf{Layer-wise representational probing}: We extract hidden states from
          every layer of \gpttwo on tool-selection prompts and train linear classifiers
          to probe (a) tool category and (b) which tool position is semantically correct.
\end{enumerate}

\para{Key findings.}  Semantic similarity predicts 40\% of LLM selections
(vs.~25\% chance), while selection remains stable across 90\% of tool orderings---showing
content dominates position.  Mechanistically, tool category is \emph{immediately} encoded
at layer~1 with 100\% probe accuracy, but identification of the \emph{correct tool
position} rises from 35\% at layer~1 to 68\% at layer~9 and then decreases in the final
layer.  This dissociation between trivial category routing and gradual semantic matching
is the paper's central mechanistic finding.

\para{Contributions.}
\begin{itemize}[leftmargin=*,itemsep=0pt,topsep=0pt]
    \item We construct a controlled tool selection dataset with 55 scenarios across five
          categories, including positional-bias and adversarial-perturbation variants,
          enabling precise behavioral and mechanistic analyses.
    \item We establish that semantic cosine similarity is a significant but incomplete
          predictor of LLM tool selection (40\% top-1 accuracy vs.\ 25\% chance), and
          that selection stability under rotation (90\%) demonstrates content dominates
          position.
    \item We provide causal evidence that tool \emph{description} content---not tool
          name---is the primary selection driver: adversarially swapping descriptions
          changes selection in 67\% of cases.
    \item We present the first layer-wise probing analysis of tool selection, revealing a
          clean dissociation: category is encoded from layer~1 (100\%), while
          within-category tool identification emerges progressively, peaking at layer~9
          (68\%), consistent with semantic matching as a multi-layer relational
          computation.
\end{itemize}

The remainder of this paper is organized as follows.  \Secref{sec:related} reviews
related work.  \Secref{sec:method} describes datasets and experimental methods.
\Secref{sec:results} presents all results.  \Secref{sec:discussion} interprets findings
and discusses limitations.  \Secref{sec:conclusion} concludes with directions for future
work.
